\chapter{StageBridge: Context-Conditioned Residual Transport} \label{chap:methods}

\section{Problem formulation}

Let $h_i\in\mathbb{R}^{d}$ denote the latent state of epithelial receiver $i$, let $e=(s,t)$ denote an ordered disease-stage edge, let $\tau\in[0,1]$ denote normalized transition time, and let $\mathcal{C}_i$ denote the typed local context surrounding receiver $i$. StageBridge estimates a continuous transition field that maps source-stage epithelial states toward target-stage epithelial states while preserving an explicit decomposition between the receiver-intrinsic transition and the context-associated residual.

The complete field is written as

\begin{equation}
\frac{d h_i}{d\tau}
=
v_{e}^{(0)}(h_i,\tau)
+
\Delta v_{\theta}(h_i,\mathcal{C}_i,e,\tau),
\label{eq:ccrt-core}
\end{equation}

where $v_{e}^{(0)}$ is the receiver-intrinsic transition field for stage edge $e$ and $\Delta v_{\theta}$ is the context-conditioned residual transport field. The central estimand is therefore

\begin{equation}
\Delta v_{\theta}
=
v_{\mathrm{full}}-v_{\mathrm{receiver}},
\label{eq:context-residual}
\end{equation}

which measures how structured local context changes the transition expected from the epithelial receiver state alone.

\section{Regulatory mediation of the context residual}

StageBridge allows the context residual to be decomposed into an interpretable regulatory contribution and a flexible nonlinear residual:

\begin{equation}
\Delta v_{\theta}(h_i,\mathcal{C}_i,e,\tau)
=
B_e r_i
+
r_{\theta}(h_i,\mathcal{C}_i,e,\tau),
\label{eq:regulatory-residual}
\end{equation}

where $r_i$ is a low-dimensional regulatory representation, $B_e$ maps regulatory programs into the direction of the stage-specific transition, and $r_{\theta}$ captures residual context effects not explained by the regulatory bottleneck alone. This formulation permits direct comparison between context-associated signaling, inferred transcription-factor programs, and the resulting transition displacement.

The regulatory representation can include pathway activity, transcription-factor activity, regulon scores, and learned latent mediators. Stage-specific mappings $B_e$ allow the same regulatory program to have different transition consequences at different points in progression.

\section{Grammar-conditioned representation}

The CCRT formulation is organized around a common transition grammar rather than a requirement that every tissue share identical cell types or molecular features. Each training instance contains a receiver state, ordered transition edge, typed sender context, relative spatial structure, signaling or pathway features, regulatory mediators, and a target transition behavior. The shared grammar can be summarized as

\begin{equation}
\mathcal{G}_i
=
\left(
\mathrm{receiver},
\mathrm{edge},
\mathrm{sender\ context},
\mathrm{distance},
\mathrm{signal},
\mathrm{regulation},
\mathrm{transition}
\right)_i.
\end{equation}

This organization provides a route to future cross-tissue evaluation without forcing LUAD and another epithelial malignancy to use identical biological vocabularies.

\section{Dual-reference latent geometry}

Each epithelial receiver and context cell is represented relative to both the Human Lung Cell Atlas and Lung Cancer Atlas. Let $z_i^{H}$ and $z_i^{L}$ denote the two reference-derived embeddings. A learned fusion module produces

\begin{equation}
\tilde{h}_i
=
f_{\mathrm{fuse}}\left(z_i^{H},z_i^{L}\right).
\end{equation}

The fusion module can assign state-dependent weight to healthy-lung and tumor-aware coordinates. Single-reference and alternative-fusion ablations test whether this learned combination improves transition modeling.

\section{Receiver-centered context tokenization}

For receiver $i$, each neighboring sender $j$ is converted into a context token

\begin{equation}
q_{ij}
=
f_{\mathrm{token}}
\left(
\tilde{h}_j,
\mathrm{type}_j,
\delta_{ij},
\pi_j
\right),
\end{equation}

where $\tilde{h}_j$ is the sender representation, $\mathrm{type}_j$ is the sender identity or compartment, $\delta_{ij}$ is the relative spatial distance, and $\pi_j$ contains pathway, signaling, or regulatory features. The resulting context is the unordered set

\begin{equation}
Q_i=\left\{q_{ij}\right\}_{j=1}^{K_i}.
\end{equation}

A receiver-centered Set Transformer encodes $Q_i$ using permutation-invariant attention. Distance enters the attention calculation as a learned bias or weight, allowing nearby senders to contribute differently from more distant senders without imposing a hard universal radius.

The context representation is

\begin{equation}
c_i
=
f_{\mathrm{context}}\left(Q_i\mid \tilde{h}_i\right),
\end{equation}

where conditioning on $\tilde{h}_i$ allows the same sender configuration to have different meanings for different epithelial receiver states.

\section{Receiver-intrinsic and full transition models}

The receiver-intrinsic model predicts the transition field using the receiver representation, stage edge, and transition time:

\begin{equation}
\hat{v}_{i}^{(0)}
=
f_{0}\left(h_{i,\tau},\tilde{h}_i,e,\tau\right).
\end{equation}

The full model adds the context representation:

\begin{equation}
\hat{v}_{i}^{\mathrm{full}}
=
f_{\mathrm{full}}
\left(h_{i,\tau},\tilde{h}_i,c_i,e,\tau\right).
\end{equation}

The learned context residual is evaluated directly as

\begin{equation}
\widehat{\Delta v}_i
=
\hat{v}_{i}^{\mathrm{full}}
-
\hat{v}_{i}^{(0)}.
\end{equation}

This paired construction is preferable to interpreting attention weights alone because it gives context a transition-level effect size in the same latent geometry as the modeled displacement.

\section{Optimal-transport coupling}

For stage edge $e=(s,t)$, let $\mu_s$ and $\mu_t$ denote empirical epithelial-state distributions at the source and target stages. StageBridge estimates an entropy-regularized optimal-transport plan

\begin{equation}
\pi_e^{*}
=
\operatorname*{arg\,min}_{\pi\in\Pi(\mu_s,\mu_t)}
\left\langle C_e,\pi\right\rangle
+
\varepsilon\,\mathrm{KL}
\left(
\pi\,\|\,\mu_s\otimes\mu_t
\right),
\label{eq:ot}
\end{equation}

where $C_e$ is a stage-specific transport cost and $\varepsilon$ controls regularization. Coupled source--target pairs are sampled from $\pi_e^{*}$ for conditional flow-matching training.

Transport costs are defined in the shared latent space and can incorporate state similarity, reference geometry, and selected biological constraints. Couplings are estimated independently for each adjacent stage edge.

\section{Conditional flow matching}

Given a coupled source state $h_0$ and target state $h_1$, an interpolation path is defined by

\begin{equation}
h_{\tau}
=
(1-\tau)h_0+\tau h_1+\sigma(\tau)\epsilon,
\end{equation}

where $\epsilon$ is a noise term and $\sigma(\tau)$ specifies the interpolation noise schedule. The conditional target velocity $u_{\tau}$ is determined by the selected probability path. The flow-matching objective is

\begin{equation}
\mathcal{L}_{\mathrm{CFM}}
=
\mathbb{E}
\left[
\left\|
\hat{v}_{\theta}(h_{\tau},\mathcal{G}_i,\tau)
-
u_{\tau}(h_0,h_1,\epsilon)
\right\|_2^2
\right].
\label{eq:cfm}
\end{equation}

The trained vector field defines a continuous stage-conditioned transition rather than only a one-step target-state prediction.

\section{Self-supervised context pretraining}

Before transition training, the context encoder is pretrained through masked receiver reconstruction. Receiver-associated molecular or state features are masked, and the model predicts the missing receiver representation from the remaining receiver information and local context. The pretraining objective is

\begin{equation}
\mathcal{L}_{\mathrm{SSL}}
=
\left\|
\hat{h}_i-h_i
\right\|_2^2,
\end{equation}

with masking and corruption applied according to the training configuration. This stage encourages the context encoder to learn biologically coherent neighborhood structure before it is optimized for transport.

\section{Growth-aware transition rates}

State displacement and population expansion are related but distinct. A growth-aware extension models the transition-associated rate as

\begin{equation}
g_i^{\mathrm{full}}
=
g_e^{(0)}(h_i,\tau)
+
\Delta g_{\phi}(h_i,\mathcal{C}_i,e,\tau),
\end{equation}

where $g_e^{(0)}$ is the receiver-intrinsic growth or survival component and $\Delta g_{\phi}$ is the context-associated residual. This component connects proliferation, survival, and tissue remodeling to the transport field without requiring them to be encoded solely as displacement.

% TODO: Align this section exactly with the frozen implementation used for the thesis.

\section{Counterfactual context perturbation}

The residual formulation supports context counterfactuals. For a selected sender population, pathway, or signaling program, a perturbed context $\mathcal{C}_i^{(-k)}$ is constructed by removal, masking, replacement, or controlled attenuation. The counterfactual effect is

\begin{equation}
\Delta_{k}v_i
=
\hat{v}^{\mathrm{full}}
\left(h_i,\mathcal{C}_i,e,\tau\right)
-
\hat{v}^{\mathrm{full}}
\left(h_i,\mathcal{C}_i^{(-k)},e,\tau\right).
\end{equation}

Counterfactual effects are summarized by receiver state, stage edge, sender type, and biological program. These analyses prioritize transition sensitivity rather than treating the presence of a ligand--receptor pair as sufficient evidence of functional influence.

\section{Training and model selection}

Model development uses patient-aware cross-validation with repeated random seeds. Hyperparameters are selected within training data, and held-out donors are used only for final fold-level evaluation. Checkpoints, model configuration, dataset manifest, random seed, fold assignment, and output schema are stored together to make each result traceable to an exact analysis run.

The total training objective combines the transition, regulatory, reconstruction, and optional growth components:

\begin{equation}
\mathcal{L}
=
\lambda_{\mathrm{CFM}}\mathcal{L}_{\mathrm{CFM}}
+
\lambda_{\mathrm{SSL}}\mathcal{L}_{\mathrm{SSL}}
+
\lambda_{\mathrm{reg}}\mathcal{L}_{\mathrm{reg}}
+
\lambda_{\mathrm{growth}}\mathcal{L}_{\mathrm{growth}}.
\end{equation}

\section{Evaluation metrics and ablations}

Primary evaluation metrics include held-out transition loss, displacement error, drift alignment, distributional distance between predicted and target states, and consistency across folds and seeds. Biological alignment is evaluated using stage-associated pathway, inflammatory, regulatory, and proliferative programs.

The core ablations are:

\begin{itemize}
    \item receiver-intrinsic model without structured context;
    \item disrupted or permuted context;
    \item context without distance information;
    \item alternative pooling, DeepSets, Set Transformer, and graph-based encoders;
    \item single-reference models using HLCA or LuCA alone;
    \item alternative dual-reference fusion strategies;
    \item removal of the regulatory bottleneck or learned gating components.
\end{itemize}

Together, these experiments test whether performance depends on biologically structured context rather than only increased model capacity.

% TODO: Add architecture diagram, algorithm environment, final hyperparameter table,
% and exact equations from the frozen CCRT implementation.
